\documentclass[11pt,a4paper]{article}

% Pakete
\usepackage[utf8]{inputenc}
\usepackage[T1]{fontenc}
\usepackage[ngerman]{babel}
\usepackage[left=18mm,right=18mm,top=18mm,bottom=18mm]{geometry}
\usepackage{helvet}
\renewcommand{\familydefault}{\sfdefault}
\usepackage{amsmath,amssymb}
\usepackage{graphicx}
\usepackage{tikz}
\usetikzlibrary{arrows.meta,positioning}
\usepackage{tcolorbox}
\tcbuselibrary{skins,breakable}
\usepackage{enumitem}
\usepackage{tabularx}
\usepackage{colortbl}
\usepackage{qrcode}
\usepackage{xcolor}
\usepackage[hidelinks]{hyperref}

% Fachfarbe Physik
\definecolor{physik}{HTML}{2c5aa0}
\definecolor{physiklight}{HTML}{e8f0fa}

% Box-Stile
\tcbset{
    kasten/.style={
        colback=white,
        colframe=physik,
        coltitle=black,
        colbacktitle=white,
        fonttitle=\bfseries,
        title=#1,
        boxrule=1pt,
        arc=2pt,
        left=5pt,
        right=5pt,
        top=5pt,
        bottom=5pt
    },
    formelbox/.style={
        colback=white,
        colframe=physik,
        boxrule=2pt,
        arc=2pt,
        left=8pt,
        right=8pt,
        top=8pt,
        bottom=8pt
    }
}

% Kopfzeile
\usepackage{fancyhdr}
\pagestyle{fancy}
\fancyhf{}
\lhead{\small Physik 12 GK}
\chead{\small Bohr-Atommodell}
\rhead{\small Name: \underline{\hspace{4cm}}}
\renewcommand{\headrulewidth}{0.4pt}

% Keine Einrueckung
\setlength{\parindent}{0pt}
\setlength{\parskip}{6pt}

\begin{document}

% Titel
\begin{minipage}[t]{0.65\textwidth}
\vspace{0pt}
{\Large\textbf{Bohr-Atommodell}}\\[0.3em]
{\normalsize Mi 28.01.2026 | Physik 12 GK}
\end{minipage}
\hfill
\begin{minipage}[t]{0.32\textwidth}
\vspace{0pt}
\raggedleft
\begin{tabular}{@{}c@{}}
\href{https://devbox42.github.io/sciencesim/physik/kl12-GK/01-photoeffekt/06-ES-bohr-postulate/LP-06-bohr-postulate.html}{\qrcode[height=1.1cm]{https://devbox42.github.io/sciencesim/physik/kl12-GK/01-photoeffekt/06-ES-bohr-postulate/LP-06-bohr-postulate.html}} \\
{\scriptsize Lernpfad}
\end{tabular}
\end{minipage}

\vspace{0.3em}
\hrule height 1pt
\vspace{0.6em}

% ============================================================
% AUFGABE 1: Das Problem
% ============================================================

\textbf{Aufgabe 1: Das Problem des Rutherford-Modells} \hfill (3 BE)

\vspace{0.3em}

a) Beobachte die Simulation im Lernpfad. Beschreibe, was mit dem Elektron passiert und warum. (2 BE)

\vspace{2cm}

b) In welcher Größenordnung liegt die vorhergesagte Zeit bis zum Kollaps? (1 BE)

\vspace{0.8cm}

% ============================================================
% AUFGABE 2: Spektrum
% ============================================================

\textbf{Aufgabe 2: Wasserstoff-Spektrum} \hfill (3 BE)

\vspace{0.3em}

a) Zeichne das Wasserstoff-Spektrum (Balmer-Serie) in den Balken ein:

\vspace{0.3em}

\begin{center}
\begin{tikzpicture}
    \draw[black, thick] (0,0) rectangle (12,0.8);
    \node[below] at (0,0) {\scriptsize 400 nm};
    \node[below] at (4,0) {\scriptsize 500 nm};
    \node[below] at (8,0) {\scriptsize 600 nm};
    \node[below] at (12,0) {\scriptsize 700 nm};
\end{tikzpicture}
\end{center}

\vspace{0.3em}

b) Das Spektrum ist \hspace{0.3cm} $\square$ kontinuierlich \hspace{0.5cm} $\square$ diskret

\vspace{0.3em}

c) Erkläre, was dieses Spektrum über die Energie im Atom aussagt. (2 BE)

\vspace{1.5cm}

% ============================================================
% FORMELKASTEN
% ============================================================

\begin{tcolorbox}[formelbox]
\textbf{Energieniveaus im Wasserstoffatom (Bohr-Modell)}

\vspace{0.5em}

\begin{center}
{\Large $E_n = \dfrac{-13{,}6 \text{ eV}}{n^2}$} \hspace{2cm} $n = 1, 2, 3, ...$
\end{center}

\vspace{0.3em}

\textbf{Photonenenergie bei Übergang:} \hspace{1cm} $E_{\text{Photon}} = |E_{n_1} - E_{n_2}| = h \cdot f$

\vspace{0.3em}

\textbf{Wellenlänge:} \hspace{1cm} $\lambda = \dfrac{h \cdot c}{E_{\text{Photon}}}$ \hspace{1cm} mit $h \cdot c = 1240 \text{ eV} \cdot \text{nm}$
\end{tcolorbox}

\vspace{0.5em}

% ============================================================
% AUFGABE 3: Berechnungen
% ============================================================

\textbf{Aufgabe 3: Energieniveaus berechnen} \hfill (4 BE)

\vspace{0.3em}

a) Berechne die Energien $E_1$, $E_2$ und $E_3$. Trage die Werte in die Tabelle ein. (2 BE)

\vspace{0.3em}

\begin{center}
\begin{tabular}{|c|c|c|c|c|c|}
\hline
$n$ & 1 & 2 & 3 & 4 & $\infty$ \\
\hline
$E_n$ (eV) & & & & $-0{,}85$ & 0 \\
\hline
\end{tabular}
\end{center}

\vspace{0.5em}

Nebenrechnung:

\vspace{1.8cm}

b) Berechne die Energie des Photons beim Übergang $n = 3 \to n = 2$. (2 BE)

\vspace{2cm}

% ============================================================
% AUFGABE 4: Anwendung
% ============================================================

\textbf{Aufgabe 4: Wellenlänge berechnen} \hfill (3 BE)

\vspace{0.3em}

Berechne die Wellenlänge des Photons aus Aufgabe 3b. In welchem Farbbereich liegt diese Linie?

\vspace{3cm}

% ============================================================
% AUFGABE 5: Transfer
% ============================================================

\textbf{Aufgabe 5: Ionisierungsenergie} \hfill (2 BE)

\vspace{0.3em}

Wie viel Energie ist nötig, um das Elektron aus dem Grundzustand ($n=1$) vollständig vom Atom zu entfernen? Begründe.

\vspace{2cm}

\vfill

\begin{center}
\rule{10cm}{0.5pt}

\textbf{Gesamt:} \underline{\hspace{1cm}} / 15 BE
\end{center}

\end{document}
